%% Bitcoin VLA Framework -- Walk-Forward 6-Fold Evaluation Report
%% Compiled on Overleaf: paste this into a new project (main.tex)
\documentclass[11pt,a4paper]{article}

\usepackage[margin=2.5cm]{geometry}
\usepackage{booktabs}
\usepackage{multirow}
\usepackage{makecell}
\usepackage[table]{xcolor}
\usepackage{array}
\usepackage{caption}
\usepackage{subcaption}
\usepackage{amsmath,amssymb}
\usepackage{graphicx}
\usepackage{hyperref}
\usepackage{enumitem}
\usepackage{times}

% ── colour helpers ──────────────────────────────────────────────────
\definecolor{bestblue}{RGB}{0, 90, 200}
\definecolor{lightgray}{gray}{0.93}
\newcommand{\best}[1]{\textcolor{bestblue}{\textbf{#1}}}
\newcommand{\refmark}{$^\dagger$}

% ── column types ────────────────────────────────────────────────────
\newcolumntype{C}[1]{>{\centering\arraybackslash}p{#1}}
\newcolumntype{L}[1]{>{\raggedright\arraybackslash}p{#1}}

% ─────────────────────────────────────────────────────────────────────
\title{\textbf{VLA-Finance: Applying Vision-Language-Action Models\\
to Bitcoin Price Prediction}\\[4pt]
\large Walk-Forward 6-Fold Evaluation Report}
\author{}
\date{\today}
% ─────────────────────────────────────────────────────────────────────

\begin{document}
\maketitle
\thispagestyle{empty}

% ─────────────────────────────────────────────────────────────────────
\section{Overview}
% ─────────────────────────────────────────────────────────────────────

We adapt the GR00T Vision-Language-Action (VLA) framework, originally developed for robot manipulation, to Bitcoin (BTC/USDT) candlestick price prediction.
The model receives an 80-candle input window (OHLCV + ETH auxiliary + text summary of recent price action) and autoregressively predicts the next 12 candles.
Five model variants are evaluated on a walk-forward 6-fold backtest spanning 2020--2025.

% ─────────────────────────────────────────────────────────────────────
\section{Experimental Setup}
% ─────────────────────────────────────────────────────────────────────

\subsection{Walk-Forward 6-Fold Protocol}

Each annual fold (fold $i$, year $Y_i$) is defined as:
\begin{itemize}[noitemsep]
  \item \textbf{Train}: Jan 1 -- Oct 31 of $Y_i$
  \item \textbf{Validation}: Nov 1 -- Nov 30 of $Y_i$ \quad (checkpoint selection)
  \item \textbf{Test}: Dec 1 -- Dec 31 of $Y_i$ \quad (held-out evaluation)
\end{itemize}
covering folds 2020--2025.
The best checkpoint per fold is selected by $\text{DA}_{k=6}$ on the validation set (tie-broken by $\text{AER}_{k=6}$).

\subsection{Trading Simulation}

An overlapping $k$-slot backtest is used: capital is split equally across $k$ independent slots, each entering/exiting every $k$ steps.
Position rule:
\[
  \text{pos}_t =
  \begin{cases}
    +1 & \hat{r}_{k} > \epsilon \\
    -1 & \hat{r}_{k} < -\epsilon \\
    0  & |\hat{r}_{k}| \leq \epsilon
  \end{cases}
\]
where $\hat{r}_{k}$ is the predicted $k$-step return and $\epsilon = 0.1\%$ is the slippage threshold.
Round-trip slippage of $0.1\%$ is deducted per trade. The benchmark is BTC buy-and-hold over the same period.

\subsection{Metrics}

\begin{description}[noitemsep]
  \item[\textbf{DA}] Directional Accuracy -- fraction of samples where predicted and actual $k$-step return share the same sign (all samples, no slippage threshold).
  \item[\textbf{Long\_ns}] Long ratio (no-slippage) -- $n_{\text{long}} / (n_{\text{long}} + n_{\text{short}})$, pure direction bias excluding watch.
  \item[\textbf{AER}] Annualised Excess Return (monthly cumulative) -- strategy return minus BTC buy-and-hold return over the test month.
  \item[\textbf{MDD}] Maximum Drawdown of the strategy equity curve.
  \item[\textbf{IC}] Trajectory Information Coefficient -- mean Pearson correlation between predicted and actual close price series (12 steps, per sample).
  \item[\textbf{RankIC}] Same as IC but using Spearman rank correlation.
\end{description}

All 6-fold results are reported as mean $\pm$ std across folds.
\refmark~v9\_delta uses a single held-out test period (2026-03) and is listed for reference only.

% ─────────────────────────────────────────────────────────────────────
\section{Models}
% ─────────────────────────────────────────────────────────────────────

\begin{description}[noitemsep]
  \item[\textbf{Kronos}~(baseline)] Off-the-shelf \texttt{NeoQuasar/Kronos-base} language model; 5 Monte-Carlo samples per prediction; no fine-tuning; 1-hour candles.
  \item[\textbf{v9\_delta}\refmark] Qwen3-VL-4B + MLP action head; delta target representation; text-only input (no chart image); 1-hour candles; single test period (reference only, not walk-forward).
  \item[\textbf{v12 GR00T}] Qwen3-VL-4B + DiT-B diffusion head (flow matching); delta target; DCT loss ($w=0.1$) + ETH auxiliary loss ($w=0.5$); \textbf{15-minute} candles.
  \item[\textbf{v13 GR00T}] Same as v12; \textbf{30-minute} candles.
  \item[\textbf{v14 GR00T}] Same as v12; \textbf{1-hour} candles.
\end{description}

% ─────────────────────────────────────────────────────────────────────
\section{Main Results}
% ─────────────────────────────────────────────────────────────────────

\subsection{Primary Comparison at \texorpdfstring{$k = 6$}{k=6}}

Table~\ref{tab:main} shows all models at prediction horizon $k=6$ (the primary evaluation horizon used for checkpoint selection).
Blue bold values indicate the best result per metric (excluding the reference-only v9\_delta).

\begin{table}[h]
\centering
\caption{6-fold average results at $k=6$. AER and MDD are cumulative over the test month.
\refmark~v9\_delta: single test period, listed for reference only.}
\label{tab:main}
\setlength{\tabcolsep}{5pt}
\renewcommand{\arraystretch}{1.25}
\begin{tabular}{lc ccc cc cc}
\toprule
\multirow{2}{*}{\textbf{Model}} & \multirow{2}{*}{\textbf{Candle}} &
  \multicolumn{2}{c}{\textbf{Direction}} &
  \multicolumn{2}{c}{\textbf{Return}} &
  \multicolumn{2}{c}{\textbf{Trajectory IC}} \\
\cmidrule(lr){3-4}\cmidrule(lr){5-6}\cmidrule(lr){7-8}
 & & DA $\uparrow$ & Long\_ns & AER $\uparrow$ & MDD $\uparrow$ & IC $\uparrow$ & RankIC $\uparrow$ \\
\midrule
\rowcolor{lightgray}
Kronos        & 1h  & $0.515_{\pm.030}$ & $0.300_{\pm.028}$ & $-0.116_{\pm.278}$ & $-0.104_{\pm.078}$ & $0.019_{\pm.065}$ & $0.015_{\pm.061}$ \\
v9\_delta\refmark & 1h & $0.544$           & $0.375$           & $+0.058$           & $-0.046$           & $-0.035$          & $-0.036$ \\
v12 GR00T     & 15m & $0.511_{\pm.022}$ & $0.442_{\pm.341}$ & $-0.374_{\pm.154}$ & $-0.349_{\pm.082}$ & $0.017_{\pm.032}$ & $0.016_{\pm.031}$ \\
v13 GR00T     & 30m & $0.513_{\pm.021}$ & $0.526_{\pm.335}$ & $-0.201_{\pm.148}$ & $-0.185_{\pm.052}$ & $0.013_{\pm.023}$ & $0.012_{\pm.023}$ \\
v14 GR00T     & 1h  & \best{$0.521_{\pm.029}$} & \best{$0.635_{\pm.297}$} & \best{$-0.093_{\pm.054}$} & $-0.124_{\pm.045}$ & \best{$0.037_{\pm.059}$} & \best{$0.039_{\pm.059}$} \\
\bottomrule
\end{tabular}
\end{table}

\subsection{Per-Horizon Breakdown}

Tables~\ref{tab:da}, \ref{tab:aer}, and \ref{tab:lns} report DA, AER, and Long\_ns across all prediction horizons $k \in \{1, 3, 6, 12\}$.

\begin{table}[h]
\centering
\caption{Directional Accuracy (DA) across horizons.}
\label{tab:da}
\renewcommand{\arraystretch}{1.2}
\begin{tabular}{lcccc}
\toprule
\textbf{Model} & $k=1$ & $k=3$ & $k=6$ & $k=12$ \\
\midrule
\rowcolor{lightgray}
Kronos            & $0.519_{\pm.011}$ & $0.524_{\pm.016}$ & $0.515_{\pm.030}$ & $0.525_{\pm.052}$ \\
v9\_delta\refmark & $0.521$           & $0.525$           & $0.544$           & $0.458$           \\
v12 GR00T         & $0.503_{\pm.010}$ & $0.504_{\pm.011}$ & $0.511_{\pm.022}$ & $0.515_{\pm.027}$ \\
v13 GR00T         & $0.498_{\pm.009}$ & $0.505_{\pm.015}$ & $0.513_{\pm.021}$ & $0.519_{\pm.023}$ \\
v14 GR00T         & \best{$0.509_{\pm.027}$} & \best{$0.525_{\pm.027}$} & \best{$0.521_{\pm.029}$} & \best{$0.527_{\pm.048}$} \\
\bottomrule
\end{tabular}
\end{table}

\begin{table}[h]
\centering
\caption{Annualised Excess Return (AER) across horizons. Higher is better.}
\label{tab:aer}
\renewcommand{\arraystretch}{1.2}
\begin{tabular}{lcccc}
\toprule
\textbf{Model} & $k=1$ & $k=3$ & $k=6$ & $k=12$ \\
\midrule
\rowcolor{lightgray}
Kronos            & $-0.334_{\pm.239}$ & $-0.176_{\pm.264}$ & $-0.116_{\pm.278}$ & $-0.100_{\pm.341}$ \\
v9\_delta\refmark & $-0.023$           & $+0.033$           & $+0.058$           & $-0.030$           \\
v12 GR00T         & $-0.998_{\pm.220}$ & $-0.664_{\pm.219}$ & $-0.374_{\pm.154}$ & $-0.206_{\pm.166}$ \\
v13 GR00T         & $-0.788_{\pm.204}$ & $-0.375_{\pm.196}$ & $-0.201_{\pm.148}$ & $-0.105_{\pm.125}$ \\
v14 GR00T         & \best{$-0.528_{\pm.154}$} & \best{$-0.174_{\pm.100}$} & \best{$-0.093_{\pm.054}$} & \best{$-0.039_{\pm.047}$} \\
\bottomrule
\end{tabular}
\end{table}

\begin{table}[h]
\centering
\caption{Long ratio (no-slippage, Long\_ns). A value of 0.5 indicates balanced directional predictions.}
\label{tab:lns}
\renewcommand{\arraystretch}{1.2}
\begin{tabular}{lcccc}
\toprule
\textbf{Model} & $k=1$ & $k=3$ & $k=6$ & $k=12$ \\
\midrule
\rowcolor{lightgray}
Kronos            & $0.406_{\pm.022}$ & $0.355_{\pm.016}$ & $0.300_{\pm.028}$ & $0.264_{\pm.031}$ \\
v9\_delta\refmark & $0.063$           & $0.103$           & $0.375$           & $0.790$           \\
v12 GR00T         & $0.411_{\pm.248}$ & $0.420_{\pm.307}$ & $0.442_{\pm.341}$ & $0.451_{\pm.362}$ \\
v13 GR00T         & $0.483_{\pm.200}$ & $0.506_{\pm.284}$ & $0.526_{\pm.335}$ & $0.544_{\pm.380}$ \\
v14 GR00T         & $0.589_{\pm.178}$ & $0.621_{\pm.252}$ & $0.635_{\pm.297}$ & $0.630_{\pm.332}$ \\
\bottomrule
\end{tabular}
\end{table}

% ─────────────────────────────────────────────────────────────────────
\section{Discussion}
% ─────────────────────────────────────────────────────────────────────

\paragraph{Candle granularity and slippage.}
The most striking trend is the monotonic improvement as candle size increases (15m $\to$ 30m $\to$ 1h).
With a fixed slippage of 0.1\% per round-trip, shorter candles incur far more frequent trades,
causing accumulated slippage to dominate and erode returns.
For example, at $k=6$ the AER improves from $-0.374$ (15m) to $-0.201$ (30m) to $-0.093$ (1h).

\paragraph{Directional bias (Long\_ns).}
v12/v13/v14 all show high variance in Long\_ns across folds (std $>$ 0.3),
indicating that the GR00T models tend to memorise the \emph{market regime} of each training period
(entirely bullish or bearish) rather than learning a generalizable directional signal.
In contrast, Kronos exhibits a systematic short bias (Long\_ns $\approx$ 0.30 at $k=6$),
which likely reflects its pre-training prior.
v14 shows the most balanced bias on average ($0.635$) but with similarly high variance.

\paragraph{Trajectory IC.}
Price IC and RankIC are near zero for all models, confirming that the models do not reliably
predict the \emph{shape} of the 12-step price trajectory beyond random chance.
v14 achieves the highest IC (0.037 $\pm$ 0.059), slightly above Kronos.

\paragraph{v14 GR00T as best model.}
At $k \in \{6, 12\}$, v14 GR00T consistently outperforms all baselines on DA, AER, IC, and RankIC
(excluding the reference-only v9\_delta).
The remaining negative AER, however, highlights that further work is needed to close the gap
to profitability.

% ─────────────────────────────────────────────────────────────────────
\section*{Appendix: Model Architecture Summary}
% ─────────────────────────────────────────────────────────────────────

\begin{table}[h]
\centering
\caption{Model configuration summary.}
\renewcommand{\arraystretch}{1.2}
\begin{tabular}{lccccc}
\toprule
\textbf{Model} & \textbf{VLM} & \textbf{Action Head} & \textbf{Target} & \textbf{Aux Loss} & \textbf{Image} \\
\midrule
Kronos        & Kronos-base & LM (direct)   & absolute  & ---       & Yes \\
v9\_delta     & Qwen3-VL-4B & MLP-2560      & delta     & ETH + DCT & No  \\
v12 GR00T     & Qwen3-VL-4B & DiT-B (flow)  & delta     & ETH + DCT & No  \\
v13 GR00T     & Qwen3-VL-4B & DiT-B (flow)  & delta     & ETH + DCT & No  \\
v14 GR00T     & Qwen3-VL-4B & DiT-B (flow)  & delta     & ETH + DCT & No  \\
\bottomrule
\end{tabular}
\end{table}

\end{document}
